\documentclass[11pt,letterpaper]{article}
\usepackage{amsmath,amssymb,amsthm}
\newtheorem{theorem}{Theorem}
\newtheorem{lemma}[theorem]{Lemma}
\theoremstyle{definition}
\newtheorem{definition}{Definition}
\title{A Note on Uniform Convergence}
\author{Marta Lindqvist}
\date{March 2024}
\begin{document}
\maketitle
\begin{abstract}
We record a short proof that a uniformly convergent sequence of continuous functions on a compact interval has a continuous limit, and we discuss a standard counterexample when uniformity fails.
\end{abstract}

\section{Definitions}
Let $f_n\colon[0,1]\to\mathbb{R}$ be a sequence of functions.

\begin{definition}\label{def:unif}
We say that $f_n\to f$ \emph{uniformly} if for every $\varepsilon>0$ there is an $N$ such that $|f_n(x)-f(x)|<\varepsilon$ for all $n\ge N$ and all $x\in[0,1]$.
\end{definition}

Pointwise convergence only asks that $N$ may depend on $x$. The difference matters, as the next section shows.

\section{The main result}
\begin{theorem}\label{thm:cont}
If each $f_n$ is continuous and $f_n\to f$ uniformly, then $f$ is continuous.
\end{theorem}
\begin{proof}
Fix $x_0$ and $\varepsilon>0$. Choose $n$ with $\sup_x|f_n(x)-f(x)|<\varepsilon/3$. By continuity of $f_n$ there is $\delta>0$ such that $|f_n(x)-f_n(x_0)|<\varepsilon/3$ whenever $|x-x_0|<\delta$. Then
\begin{align}
|f(x)-f(x_0)| &\le |f(x)-f_n(x)|+|f_n(x)-f_n(x_0)|+|f_n(x_0)-f(x_0)| \nonumber\\
&< \varepsilon. \label{eq:three}
\end{align}
Hence $f$ is continuous at $x_0$.
\end{proof}

The estimate \eqref{eq:three} is the familiar ``$\varepsilon/3$ argument''; Theorem~\ref{thm:cont} follows from it.

\section{A counterexample}
Take $f_n(x)=x^n$. Then $f_n(x)\to0$ for $x<1$ while $f_n(1)=1$, so the limit is discontinuous. By Theorem~\ref{thm:cont} the convergence cannot be uniform; indeed $\sup_x |f_n(x)-f(x)|=1$ for every $n$.
\end{document}
